%%%% utfpr-ict-report.tex, 2024/10/05
%%%% Copyright (C) 2017-2024 Luiz E. M. Lima (luizeduardomlima@gmail.com)
%%
%% This work may be distributed and/or modified under the conditions of the
%% LaTeX Project Public License, either version 1.3 of this license or (at your
%% option) any later version.
%% The latest version of this license is in
%%   https://www.latex-project.org/lppl.txt
%% and version 1.3 or later is part of all distributions of LaTeX version
%% 2005/12/01 or later.
%%
%% This work has the LPPL maintenance status `maintained'.
%%
%% The Current Maintainer of this work is Luiz E. M. Lima.
%%
%% This work consists of all files listed in readme.html.
%%
%% A brief description of this work is in readme.txt.

%% Detecção e aviso sobre comandos obsoletos
% \RequirePackage[l2tabu, orthodox]{nag}

%% Classe de documento
\documentclass[%% Opções (^ = padrão; > = para pacotes):
  a4paper,%% Tamanho de papel: a4paper, letterpaper (^), etc.
  12pt,%% Tamanho de fonte: 10pt (^), 11pt ou 12pt
  fleqn,%% Alinhamento de equações à esquerda (centralizado por padrão)
%   draft,%% Aparência de documento (>): draft ou final (^)
  english,%% Idioma secundário (penúltimo) (>)
  brazilian,%% Idioma primário (último) (>)
]{article}

%% Pacotes utilizados
\usepackage[%% Opções (^ = padrão):
%   Font      = Times,%% Fonte principal: Arial (^), Calibri, Times ou LaTeX(*)
%   LinkColor = ,%% Cor (Modelo HTML) hyperlink (azul escuro ^): nenhum modelo (cor texto)
%   CvrHdr    = None,%% Cabeçalho na Capa: UTFPR (^) ou None
%   DocLic    = CC,%% Licença do documento: CC (Creative Commons) ou None (^)
%   CCLic     = BY-NC-ND,%% Licença CC: BY (^), BY-SA, BY-ND, BY-NC, BY-NC-SA ou BY-NC-ND
%   ABNTCit   = AAY,%% Citação ABNT: AAY (Nome, Ano), NRB (1) ou NSB [1] (^)
%   DOIIcon   = On,%% Ícone de DOI em referências: On ou Off (^)
%   URLIcon   = On,%% Ícone de URL em referências: On ou Off (^)
]{utfpr-ict-report}

%% Informações do documento (descomentar para alterar; * = opcional)
%%%% Trabalho: {Título} (conforme Plano de Trabalho)
% \Title{Título do Trabalho do{(a)} Estudante}
%%%% Projeto: {Título} (conforme Plano de Trabalho)
% \ProjectTitle{Título do Projeto do{(a)} Orientador{(a)}}
%%%% Estudante: {Dados}
% \Student{%
%   Name      = {Nome Completo do{(a)} Estudante},%
%   Email     = {student@domain},%% (*)
%   Lattes    = {0000000000000001},%% (*)
%   ORCID     = {0000-0000-0000-0001},%% (*) (CHKTEX 8)
%   Course    = {Engenharia Mecânica},%
%   Mode      = {Bolsista PIBIC/CNPq},%% Conforme Plano de Trabalho
%   Signature = {example-image},%% Arquivo de imagem (*) para submissão
% }
%%%%% Orientador(a): {Dados}
% \Advisor{%
%   Title     = {Prof{(a)}. Dr{(a)}.},%
%   Name      = {Nome Completo do{(a)} Orientador{(a)}},%
%   Email     = {advisor@domain},%% (*)
%   Lattes    = {0000000000000002},%% (*)
%   ORCID     = {0000-0000-0000-0002},%% (*) (CHKTEX 8)
%   Mode      = {Orientador{(a)}},%% Conforme Plano de Trabalho
%   Signature = {example-image},%% Arquivo de imagem (*) para submissão
% }
%%%% Coorientador(a) (se houver): {Dados}
% \CoAdvisor{%
%   Title     = {Prof{(a)}. Dr{(a)}.},%
%   Name      = {Nome Completo do{(a)} Coorientador{(a)}},%
%   Email     = {coadvisor@domain},%% (*)
%   Lattes    = {0000000000000003},%% (*)
%   ORCID     = {0000-0000-0000-0003},%% (*) (CHKTEX 8)
%   Mode      = {Coorientador{(a)}},%% Conforme Plano de Trabalho
%   Signature = {example-image},%% Arquivo de imagem (*) para submissão
% }
%%%% Relatório: [Etapa] (Final ou Parcial); {Tipo} (de Pesquisa ou Técnico)
% \ReportInfo[Final]{de Pesquisa}
%%%% Iniciação: [Ingresso] (Mês/Ano); {Tipo} (Científica ou Tecnológica)
% \InitiationInfo[08/2023]{Científica}
%%%% Área do Conhecimento (CNPq): [Código]; {Nome}
% \KnowledgeArea[3.00.00.00{-}9]{Engenharias}
%%%% Campus: {Cidade}
% \Campus{Ponta Grossa}
%%%% Data: {Ano} (atual por padrão)
% \Date{2024}
%%%% Atribuição de Licença CC (*): {Ano}; {Autoria}
% \CCLicenseAttribution{\YearNum}{N. C. Estudante {\&} N. C. Orientador{(a)}}

%% Resumo e Palavras-chave
\begin{abstract}[After]%% Localização relativa ao sumário: [Before], [With] ou [After]
Deve ser redigido de forma clara e sucinta na terceira pessoa do singular, com verbo na voz ativa, contendo de 150 a 200 palavras e seguido de três a cinco palavras-chave.
Deve situar o trabalho no contexto geral e apresentar objetivos, metodologia adotada, contribuição própria, resultados obtidos e conclusões mais importantes do trabalho.
\keywords{palavra-chave-1\and palavra-chave-2\and palavra-chave-3}
\end{abstract}

%% Legendas de figuras e tabelas: {Estilo}
% \captionsetup{font+ = bf}%% Negrito (texto normal por padrão)

%% Arquivo(s) de referências
\addbibresource{utfpr-ict-report.bib}

%% Início do documento
\begin{document}

\section{Introdução}%
\label{sect:intro}

Deve apresentar uma visão global da pesquisa, incluindo: breve histórico, importância e justificativa de escolha do tema, delimitações do assunto, formulação de hipóteses e objetivos da pesquisa.

\UTFPR-ICT-Report é um modelo \LaTeX\footnote{Para compilação, usar \texttt{pdflatex} (preferencialmente), \texttt{lualatex} ou \texttt{xelatex} (em último caso)} e tem por finalidade padronizar os relatórios de atividades de Iniciação Científica ou Tecnológica da \href{https://www.utfpr.edu.br/}{UTFPR\LinkIcon}, realizadas pelos{(as)} orientados{(as)} dos programas PIBIC, PIBIC-Af, PIBIC-EM, PIBIT, PIVIC e PIVIT\@.
Relatórios submetidos podem ser rejeitados caso não sigam este padrão.

Os relatórios assinados digitalmente por professor{(a)}-orientador{(a)} e estudante, necessariamente em formato PDF\footnote{Pode ser convertido para PDF/A em: \url{https://www.pdfforge.org/online/en/pdf-to-pdfa}} (\ENLang*{Portable Document Format} ou Formato de Documento Portátil), devem ser submetidos utilizando o Sistema de Pesquisa (SISPEQ) nos \href{https://sistemas2.utfpr.edu.br/}{Sistemas Corporativos Integrados\LinkIcon} da \href{https://www.utfpr.edu.br/}{UTFPR\LinkIcon}, em data a ser definida pela \PROPPGName\ (PROPPG).

\section{Material e Métodos}%
\label{sect:mat-meth}

Deve apresentar modelos utilizados, modelagem empregada, simplificações necessárias, metodologia e descrição do método de cálculo utilizado no desenvolvimento da pesquisa, para que a mesma possa ser reconstituída.
Devem ser descritos também: montagem experimental, metodologia para a obtenção de resultados, análise de erros, amostras de resultados obtidos e comentários.
Esta seção pode ter outro título conforme o trabalho realizado.

\subsection{Tamanho do trabalho}%
\label{ssect:work-size}

O trabalho completo, incluindo figuras e tabelas, deve conter de 3 (três) a 15 (quinze) páginas, em papel de tamanho padrão A4 (\Unit[21]{cm} \texttimes\ \Unit[29,7]{cm}).
Figuras e tabelas não devem ser reduzidas a tamanhos que sacrifiquem o entendimento dos símbolos e legendas contidos nos mesmos.

\subsection{Configuração da página}%
\label{ssect:page-config}

Cada página, no tamanho A4, deve ser formatada de modo a apresentar \Unit[2,5]{cm} de margem em todos os lados do documento.
Dentro desta área, o texto deve ser formatado em uma única coluna, sem incluir moldura no texto.

\subsection{Formatação do texto}%
\label{ssect:text-format}

O texto deve ser iniciado pela \nameref{sect:intro}.
Os títulos das seções primárias (\nameref{sect:intro}, \nameref{sect:mat-meth}, etc.) devem ser escritos em negrito, sem numeração, em maiúsculo e alinhados à esquerda, sendo que o conteúdo, propriamente dito, deve ser iniciado após espaçamento de uma linha em branco e \Unit[1]{cm} de recuo.
Ao final de cada seção primária, deve-se deixar uma linha em branco.
Todo o texto deve ser escrito com espaçamento de 1,5 (entre linhas), usando fonte Arial ou Times New Roman (ou equivalentes) de tamanho \Unit[12]{pt}.
Para as seções secundárias, o título também deve ser escrito em negrito, mas somente com a primeira letra em maiúscula, além de apresentar alinhamento à esquerda seguido por um ponto.
Inicie o texto da seção secundária imediatamente após o título da mesma.
As notas de rodapé\footnote{Exemplo de nota de rodapé} devem ser colocadas na parte inferior da página correspondente (usando fonte de tamanho de \Unit[8]{pt}) e separadas do texto por um traço.

\section{Resultados e Discussão}%
\label{sect:res-disc}

Deve descrever detalhadamente os dados obtidos no trabalho.
Os resultados são normalmente discutidos a partir de ilustrações (gráficos, quadros, etc.), tabelas, entre outros elementos, que podem ser incluídos no documento.
Deve efetuar a comparação dos dados obtidos e/ou resultados com aqueles descritos na revisão da literatura, incluindo os comentários sobre os estudos de outros trabalhos.

As referências devem ser listadas no fim do relatório, na ordem de citação, conforme formato da Associação Brasileira de Normas Técnicas (ABNT).
No texto, as citações devem ser referenciadas por seu número de ocorrência colocado entre colchetes, por exemplo, as citações indiretas:

\begin{itemize}
\item Implícitas:
\begin{itemize}
\item \ldots\ \cite{Ekenstein1997}.
\item \ldots\ \cite{Coleman1991,Nriagu1988}.
\item \ldots\ \cite{Wizentier1992,Faina2000,Larsson2020}.
\item \ldots\ \cite{Ekenstein1997,Nriagu1988,Faina2000}.
\end{itemize}
\item Explícitas:
\begin{itemize}
\item \textcite{Ekenstein1997} afirmam que\ldots%
\item \ldots\ conforme \textcite{Coleman1991,Nriagu1988}.
\item Segundo \textcite{Wizentier1992,Faina2000,Larsson2020}, \ldots%
\item \ldots\ as definições de \textcite{Ekenstein1997,Nriagu1988,Faina2000}.
\end{itemize}
\end{itemize}

Citações diretas curtas (de até três linhas) acompanham o corpo do texto e se destacam com aspas duplas.
Caso o texto original já contenha aspas, estas devem ser substituídas por aspas simples.
Exemplo: \Citation[brazilian]{\cite[\ppno~20, grifo do autor]{Fulano2021}}{\textelp{} citar trechos de \enquote{outros autores} sem os referenciar, pode ser caracterizado plágio}.
Para as citações diretas longas (com mais de três linhas), estas devem ser transcritas em parágrafo distinto.
Por exemplo:

\begin{DisplayCitation}[brazilian]{\cite[150]{Fulano2021}}
Toda citação direta com mais de três linhas é considerada uma citação direta longa.
Este tipo de citação deve ser escrita sem aspas, em parágrafo distinto, com fonte de tamanho \Unit[10]{pt}, espaçamento simples e com recuo de \Unit[4]{cm} da margem esquerda, terminando na margem direita, conforme ilustrado neste exemplo
\end{DisplayCitation}

Conforme exemplos no arquivo-fonte deste modelo, citações diretas curtas e longas podem ser inseridas por meio do pacote \LaTeX\ \texttt{utfpr-ict-report} usando:

\begin{itemize}
\item Comando \LaTeX\ \verb|\Citation|.
\item Ambiente \LaTeX\ \texttt{DisplayCitation}.
\end{itemize}

Citações e referências podem ser inseridas usando os comandos do pacote \href{https://ctan.org/pkg/biblatex/}{Bib\LaTeX\LinkIcon}, conforme exemplos no arquivo-fonte deste modelo.
Os dados de cada referência podem ser obtidos de um arquivo \href{https://www.bibtex.org/}{Bib\TeX\LinkIcon} (\texttt{*.bib}), geralmente na própria página de acesso ou \ENLang{download} da publicação (artigos, livros, etc.) ou, ainda, a partir do Google Acadêmico, etc.
Diversas ferramentas podem ser usadas para gerar ou editar entradas de arquivos \href{https://www.bibtex.org/}{Bib\TeX\LinkIcon} (\texttt{*.bib}), por exemplo:

\begin{itemize}
\selectlanguage{english}
\item \href{https://zbib.org/}{ZoteroBib\LinkIcon}.
\item \href{https://truben.no/latex/bibtex/}{Bib\TeX\ Editor\LinkIcon}.
\end{itemize}

Todas as equações devem ter \Unit[1]{cm} de recuo da margem esquerda e ser numeradas sequencialmente, com os números entre parênteses na margem direita, conforme o exemplo da \cref{eq:et}:
\begin{equation}%
\label{eq:et}
e{\left(t\right)} = \sum_{n = 1}^5 \frac{1}{2 + n} \cos{\left(2 \pi n t\right)}
\end{equation}

As equações devem ser referenciadas no texto da seguinte forma: \enquote{Substituindo a \cref{eq:et} na \textelp{}, obtém-se \textelp{}}.

Equações podem ser inseridas usando o ambiente \LaTeX\ \texttt{equation}, conforme exemplo no arquivo-fonte deste modelo.
Símbolos matemáticos (ou equações mais simples) podem ser inseridos ao longo do texto de um parágrafo usando o ambiente \LaTeX\ \texttt{math} (ou o atalho \LaTeX\ \verb|\(...\)| ou ainda o atalho \TeX\ \verb|$...$|), por exemplo: $\alpha$, $A = \pi D^2 / 4$, etc.
Diversas ferramentas podem ser usadas para gerar ou editar equações em \LaTeX, por exemplo:

\begin{itemize}
\selectlanguage{english}
\item \href{https://formulasheet.com/}{Formula Sheet\LinkIcon}.
\item \href{https://www.tutorialspoint.com/latex_equation_editor.htm}{\LaTeX\ Equation Editor (\textit{by} Tutorials Point)\LinkIcon}.
\end{itemize}

Todos os elementos gráficos, como figuras e tabelas, devem ser centralizados.
As figuras e tabelas (entre outros elementos) devidamente referenciadas no texto, podem ser colocadas da maneira mais conveniente para os autores em uma ou duas colunas, desde que o texto permaneça em apenas uma coluna.
Antes e após os elementos não textuais e suas respectivas legendas, deve-se deixar uma linha de espaçamento.
Os autores não devem se esquecer da colocação de legendas nas figuras, tabelas e outros elementos gráficos inseridos no texto.

As figuras devem ser numeradas sequencialmente com algarismos arábicos, conforme os exemplos das \cref{fig:t-x,fig:campi-map}.

\begin{figure}[!htb]
\savebox0{\includegraphics[width = 125mm]{fig-campi-map}}
\usebox0%
\llap{%
  \raisebox{\ht0-\height}{%
    \href{https://www.utfpr.edu.br/campus}{%
      \includegraphics[height = 15mm]{fig-campi-map-qr-code}%
    }%
  }%
}
\caption{Localização dos campi da \textcite{UTFPR2017}.}%
\label{fig:campi-map}
\end{figure}

\begin{figure}[!htb]
\includegraphics[width = 0.8\linewidth]{grph-t-x}
\caption{Centralizada na coluna e com legenda abaixo da mesma.}%
\label{fig:t-x}
\end{figure}

Figuras podem ser inseridas usando o ambiente \LaTeX\ \texttt{figure}, conforme exemplos no arquivo-fonte deste modelo.
A \cref{fig:t-x} foi produzida usando o ambiente \LaTeX\ \texttt{tikzpicture} do pacote \LaTeX\ \texttt{tikz} a partir do arquivo \texttt{grph-t-x.tex} em \texttt{./Figures/}.
Como mostrado na \cref{fig:campi-map}, objetos flutuantes (ilustrações) podem adicionalmente receber um código QR (\ENLang*{Quick Response} ou Resposta Rápida) contendo: URL (\ENLang*{Uniform Resource Locator} ou Localizador Uniforme de Recursos) ou informações complementares.

Todas as tabelas devem ser numeradas sequencialmente com algarismos arábicos, conforme o exemplo da \cref{tab:cases-abc}.

\begin{table}[!htb]
\caption{Legenda acima da mesma, centralizada.}%
\label{tab:cases-abc}
\begin{tabular}{@{}lcccc@{}}
\toprule%
\multicolumn{1}{@{}c}{Caso} & $w$ & $x$ & $y$ & $z$ \\
                            & (m) & (m) & (m) & (m) \\ \midrule%
A                           & 1   & 4   & 7   & 10  \\
B                           & 2   & 5   & 8   & 11  \\
C$^{(1)}$                   & 3   & 6   & 9   & 12  \\ \bottomrule%
\multicolumn{5}{@{}l@{}}{$^{(1)}$ Exemplo de nota.}
\end{tabular}
\end{table}

Tabelas podem ser inseridas usando o ambiente \LaTeX\ \texttt{table}, conforme exemplo no arquivo-fonte deste modelo.
Diversas ferramentas podem ser usadas para gerar ou editar tabelas em \LaTeX, por exemplo:

\begin{itemize}
\selectlanguage{english}
\item \href{https://www.tablesgenerator.com/}{Tables Generator\LinkIcon}.
\item \href{https://www.latex-tables.com/}{\LaTeX\ Tables Editor\LinkIcon}.
\end{itemize}

Ressalta-se que as legendas das figuras devem ser colocadas abaixo das mesmas, enquanto as das tabelas devem ser colocadas acima das mesmas.

Informações e dicas sobre \TeX/\LaTeX\ podem ser obtidas em:

\begin{enumerate}[label = {\alphsect}]
\selectlanguage{english}
\item \href{https://www.latex-project.org/}{\LaTeX\ Project\LinkIcon};
\item \href{https://www.ctan.org/}{Comprehensive \TeX\ Archive Network (CTAN)\LinkIcon};
\item \href{https://www.tug.org/}{\TeX\ Users Group (TUG)\LinkIcon};
\item \href{https://en.wikibooks.org/wiki/LaTeX/}{\LaTeX\ \textemdash\ Wikibooks\LinkIcon};
\item \href{https://tex.stackexchange.com/}{\TeX-\LaTeX\ Stack Exchange\LinkIcon}.
\end{enumerate}

Pede-se, ainda, que todos os autores revisem cuidadosamente a versão final do trabalho para evitar erros de digitação e formatação.
Os relatórios devem ser submetidos via formulário eletrônico, utilizando o Sistema de Pesquisa (SISPEQ) nos \href{https://sistemas2.utfpr.edu.br/}{Sistemas Corporativos Integrados\LinkIcon} da \href{https://www.utfpr.edu.br/}{UTFPR\LinkIcon}, anexando-se o relatório no formato PDF (devendo ser necessariamente \enquote{sem proteção}).
É de fundamental importância que os relatórios estejam devidamente assinados digitalmente por orientador{(a)} e estudante.

\section{Conclusões}%
\label{sect:concl}

Deve finalizar o trabalho com respostas às hipóteses especificadas na introdução.
O ponto de vista sobre os resultados obtidos deve ser expresso; não se deve incluir novos dados ou equações.
A partir da tese, alguns assuntos identificados como importantes para serem explorados podem ser sugeridos como temas para novas pesquisas.

Bolsistas e voluntários{(as)}, participantes dos programas de Iniciação Científica ou Tecnológica da \href{https://www.utfpr.edu.br/}{UTFPR\LinkIcon}, devem enviar os relatórios pelo meio recomendado até a data requisitada, para cumprir os requisitos do programa e ter direito ao certificado.

\section{Agradecimentos}%
\label{sect:ack}

Havendo agradecimentos, estes devem vir antes das \nameref{sect:ref}.
Bolsistas obrigatoriamente devem agradecer à agência de fomento da bolsa, por exemplo:

\begin{itemize}
\item O presente trabalho foi realizado com apoio do CNPq, Conselho Nacional de Desenvolvimento Científico e Tecnológico \textemdash\ Brasil \textemdash\ Edital UTFPR/PROPPG n.\textsuperscript{o} XX/YYYY (PIBIC).
\item O presente trabalho foi realizado com apoio da Fundação Araucária \textemdash\ Brasil \textemdash\ Edital UTFPR/PROPPG n.\textsuperscript{o} XX/YYYY (PIBIT).
\item O presente trabalho foi realizado com apoio da UTFPR, \UTFPRName\ \textemdash\ Brasil \textemdash\ Edital UTFPR/PROPPG n.\textsuperscript{o} XX/YYYY (PIVIC/PIVIT).
\end{itemize}

\printbibliography%% Referências

%% Fim do documento
\end{document}
